\section{Processi gestionali}

\subsection{Obiettivi e priorità}
L'obiettivo di gestione è consegnare i sette deliverable (D1--D7) nelle date fissate dal docente e arrivare a un'app Android funzionante che copra i requisiti a priorità Alta del D1. Tra i tre vincoli di progetto le priorità sono:
\begin{enumerate}
    \item \textbf{Tempi}: le scadenze dei deliverable sono fisse e non negoziabili, quindi hanno la precedenza su tutto.
    \item \textbf{Qualità}: a parità di tempo si preferisce consegnare meno funzionalità ma verificate, partendo dai requisiti RF a priorità Alta (uso offline, schede, mappa, foto e note georeferenziate, esportazione).
    \item \textbf{Costi}: il progetto non ha budget economico; il costo coincide con le ore di lavoro dei membri e con strumenti gratuiti o open source, per cui non è un vincolo critico.
\end{enumerate}
In caso di ritardo, l'ambito si riduce rinviando i requisiti a priorità Media e Bassa, senza spostare le scadenze.

\subsection{Monitoraggio e controllo}
L'avanzamento è controllato rispetto al diagramma di Gantt (Sezione~5.3) e alla ripartizione dei compiti di ogni deliverable (Sezione~2).
\begin{itemize}
    \item \textbf{Riunione settimanale} del team (in presenza o online) per verificare lo stato delle attività, i problemi aperti e i compiti della settimana successiva; le decisioni sono riportate in un breve verbale (Appendice~B).
    \item \textbf{Stato delle attività}: ogni attività è tracciata come \emph{issue} su GitHub con responsabile e stato (da fare, in corso, completata), così l'avanzamento è visibile a tutti.
    \item \textbf{Indicatori}: percentuale di attività completate rispetto al piano, numero di requisiti RF implementati e verificati, issue aperte e giorni mancanti alla scadenza del deliverable.
    \item \textbf{Revisione prima della consegna}: alcuni giorni prima di ogni scadenza il team rilegge insieme il documento compilato e verifica che tutte le sezioni siano complete.
\end{itemize}
Se un'attività accumula ritardo, il responsabile lo segnala subito sul gruppo; nella riunione successiva il team ridistribuisce il lavoro o riduce l'ambito secondo le priorità indicate sopra. Il referente di progetto, Ibrahima Tely Barry, interviene se il ritardo mette a rischio una consegna.

\subsection{Gestione dei rischi}
I rischi sono stati individuati dal team e valutati per probabilità e impatto (Alta, Media, Bassa). Sono riesaminati nella riunione settimanale e aggiornati a ogni deliverable.
\begingroup
\SetTblrTemplate{contfoot}{empty}
\SetTblrTemplate{conthead}{empty}
\DefTblrTemplate{conthead-text}{nocont}{}
\DefTblrTemplate{contfoot-text}{nocont}{}
\SetTblrTemplate{conthead-text}{nocont}
\SetTblrTemplate{contfoot-text}{nocont}
\begin{longtblr}[entry=none,label=none]{width=\linewidth,
             colspec={X[0.4,l]X[l]X[0.7,l]X[0.7,l]X[1.4,l]}, rowhead=1, cells={font=\small},
             row{1}={bg=arancio}, hline{2-Z}={wd=0.5pt, fg=arancio}}
ID & Rischio & Probabilità & Impatto & Mitigazione \\
R1 & Ritardo nella consegna di un deliverable & Media & Alta & Attività tracciate su GitHub, revisione alcuni giorni prima della scadenza, riduzione dell'ambito secondo le priorità \\
R2 & Poca esperienza del team nello sviluppo Android & Alta & Media & Prototipi semplici all'inizio, documentazione ufficiale, sviluppo per moduli successivi \\
R3 & Requisiti instabili o non validati dai docenti & Media & Alta & Validazione dei requisiti con i docenti prima dello sviluppo, modifiche registrate come issue \\
R4 & Complessità del funzionamento offline e della sincronizzazione & Alta & Alta & Architettura offline-first fin dall'inizio, test senza connessione, sincronizzazione semplice prima di quella automatica \\
R5 & Sensori (GPS, bussola, clinometro) assenti o imprecisi su alcuni dispositivi & Media & Media & Test su dispositivo fisico, inserimento manuale sempre disponibile come alternativa \\
R6 & Indisponibilità temporanea di un membro (esami, malattia) & Media & Media & Rotazione dei compiti, codice e documenti condivisi su GitHub, sostituzione da parte di un altro membro \\
R7 & Perdita di lavoro o conflitti tra versioni & Bassa & Alta & Git e GitHub, commit frequenti, pull request revisionate, nessuna modifica diretta sul ramo principale \\
R8 & Scarso coinvolgimento degli utenti finali (il questionario del D1 ha raccolto 7 risposte) & Media & Media & Nuove occasioni di confronto con docenti e studenti, prove di usabilità sul campo \\
\end{longtblr}
\endgroup

\subsection{Gestione delle comunicazioni}
\begin{itemize}
    \item \textbf{Interne al team}: un gruppo WhatsApp per il coordinamento quotidiano e le domande rapide; riunioni settimanali per decisioni e avanzamento; GitHub (issue e pull request) per tutto ciò che riguarda codice e documenti, così le discussioni tecniche restano legate al lavoro a cui si riferiscono.
    \item \textbf{Con il docente}: solo tramite il referente di progetto, che invia le consegne su Moodle e gestisce le comunicazioni formali (Sezione~2.3).
    \item \textbf{Con gli utenti finali}: questionari (Google Form) e incontri, anche online su Zoom, con i docenti e gli studenti del Laboratorio in Campo per raccogliere e validare i requisiti.
\end{itemize}
Le decisioni importanti prese in chat vengono riportate nel verbale della riunione successiva, in modo che restino documentate.

\subsection{Gestione della configurazione e della qualità}
Codice e documenti sono gestiti con Git in un repository GitHub condiviso. Le modifiche entrano nel ramo principale tramite pull request revisionate da un altro membro, e ogni consegna è contrassegnata da un tag che permette di ritrovare la versione consegnata. Ogni documento e ogni funzione sono controllati da un membro diverso dall'autore prima della consegna o dell'integrazione; un'attività è considerata conclusa solo dopo questo controllo. Le regole operative di versionamento e il piano della qualità (chi controlla, che cosa e quando) sono descritti nella Sezione~4.3.
